%% =================================================================
%% Chen et al. Measurement 264 (2026) 120300.
%% Reconstruction of page 5 (printed p. 5).
%% Contents: Fig. 7 (a-d: hierarchical thickness study with 0.5 mm
%% step model and compression-depth gradation), plus Sections 2.2
%% (Probe module) and 2.3 (Closed-loop feedback control system) and
%% the start of Section 3 / 3.1 (Softness recognition).
%% No numbered equations on this page.
%% =================================================================

\documentclass[11pt]{article}

\usepackage[margin=0.85in]{geometry}
\usepackage{amsmath, amssymb}
\usepackage{mathrsfs}
\usepackage{xcolor}
\usepackage{fancyhdr}

\pagestyle{fancy}
\fancyhf{}
\lhead{\small Z.~Chen et al.}
\rhead{\small Measurement 264 (2026) 120300}
\cfoot{\small 5 $\to$ \arabic{page}}
\renewcommand{\headrulewidth}{0.4pt}

\setcounter{page}{1}
\setlength{\parskip}{6pt}
\setlength{\parindent}{0pt}

\begin{document}

% ---------- Fig. 7 placeholder ----------
\begin{center}
\fbox{\begin{minipage}[c][6.5cm][c]{0.95\textwidth}
\centering
\textbf{Fig. 7.}\ Effect of the thickness of the sample on the
tactile softness of materials.\\[4pt]
\textcolor{gray}{\footnotesize
\textbf{(a)} Hierarchical structure model with 9 staircase steps,
0.5 mm height difference between adjacent stages (total run 5 mm),
CAD view with the stepped top surface.\\[3pt]
\textbf{(b)} Hierarchical model filled with a soft surface layer
of AB glue (photograph), uniform square scanning region outlined.\\[3pt]
\textbf{(c)} 3D surface plot of compression-depth distribution when
the tactile tomography system scanned (b); stepped pyramid shape
with X, Y in mm and Z 0--6 mm on a rainbow color scale.\\[3pt]
\textbf{(d)} Bar chart: measured height (mm, 0--7 axis) at step
numbers 1--9; consecutive step-to-step height differences are
0.494 mm and 0.495 mm, matching the 0.5 mm nominal staircase step.]}
\end{minipage}}
\end{center}

\vspace{0.4em}

\textbf{Fig. 7.}\ Effect of the thickness of the sample on the
tactile softness of materials. (a) Hierarchical structure model
with the height difference between each adjacent stage of 0.5 mm.
(b) Hierarchical structure model filled with a soft surface layer
of AB glue. (c) 3D imaging of compression depth distribution when
the tactile tomography system scanned the hierarchical structure
model with a soft surface layer. (d) Height of the compression
depth distribution.

\vspace{0.8em}

% ---------- Prose continuation (from page 4) ----------
0.1 $\mu$m, respectively, enabling the tactile tomography system
to operate with high scanning resolution.

\subsection*{2.2.\ Probe module}

As shown in Figs.~2 and 3, the probe module of the tactile
tomography system was assembled from a shell structure, a tip, two
springs, and a carbon fiber beams based pressure sensor [35]. To
achieve high scanning resolution without penetrating the soft layer
of the sample, the tip was designed with a diameter of 0.5 mm. The
springs are positioned behind the tip to transmit the force from
the tip to the pressure sensor and to restore the tip's position.
The pressure sensor, which is based on carbon fiber beams, is
capable of quantitatively measuring force and responds to the force
exerted by the springs. Consequently, the probe module can
accurately capture force information when the tip makes contact
with the sample.

\subsection*{2.3.\ Closed-loop feedback control system}

A programmable logic controller (PLC) (H2U-48MT, Henchman) was
integrated with the scanning module and the probe module to
establish a closed-loop feedback control system. The workflow of
the tactile tomography system is shown in Fig.~4. Firstly, scanning
paths, scanning speed, and scanning step were programmed. The PLC
control the X and Y axes to drive the sample to move until the
target scanning point is right under the tip, and then the rodless
cylinder of the Z axis drive the probe down. The resistance value
of the pressure sensor changes when the tip of the probe module
touches the surface of the sample. This resistance value is
collected by an ADS1247 acquisition chip, which converts the analog
signal into a digital signal. The digital signal is then
transmitted to an STM32F030C8T6 microprocessor, where it is
identified as the first threshold. Meanwhile, the position
information of $Z_0$ is read by the grating sensor, and the
microprocessor sends this position information to the PLC. The
rodless cylinder then continues to drive the probe downward, while
the grating sensor reads a series of positional data as the
resistance value of the pressure sensor reaches specific
characteristic thresholds. The PLC stops and resets the probe once
the resistance value of the pressure sensor attains the target
threshold. The X and Y axes then drive the sample to the next
scanning point until all preset areas have been scanned. The
positional information $(x, y, z)$ corresponding to each threshold
is saved as projection data to reconstruct a 3D image of the
internal structure of the sample. This projection data is stored
in a CSV file and displayed in real time through a visual
interface.

\section*{3.\ Imaging principle of tactile tomography}

\subsection*{3.1.\ Softness recognition of tactile tomography system}

The ability of a material sample to resist the penetration of hard
objects into its surface is known as hardness, which is an intrinsic
property of the material. [36]. According to different measurement
approach and standard, hardness can be subdivided into Rockwell
hardness, Vickers hardness, Shore hardness, and Brinell hardness.
Among these hardness measurements, Shore hardness has been widely
applied to the hardness testing of soft materials, such as
elastomers, rubber, soft plastics, and foam [37]. Therefore, the
softness of soft matter can be characterized by Shore hardness,
which exhibits a negative correlation with softness. The tactile
tomography system can

\end{document}
